%%%%%%%%%%%%%%%%%%%%%---Chapter4 Formatting---%%%%%%%%%%%%%%%%%%%%%%%%%%%%%%%%%%
\documentclass[Thesis.tex]{subfiles}
\begin{document}
\newpage
\setcounter{chapter}{3} %%Chapter Numbering

\titleformat{\chapter}[display]
{\sffamily\fontsize{27}{27}\bfseries\filleft}{\thechapter}{0pt}{{#1}}

\thispagestyle{empty}
%%%%%%%%%%%%%%%%%%%%%%%%%%%%%%%%%%%%%%%%%%%%%%%%%%%%%%%%%%%%%%%%%%%%%%%%%%%%

\chapter{Experimental Setup and Implementation}\label{chap4:setup}

This chapter records the implementation in enough detail to reproduce every number
in Chapters~\ref{chap5:results} and~\ref{chap6:decoding}. The intent is that a
reader with the accompanying source tree and these version pins can regenerate the
result tables and figures without contacting the author.

\section{Software stack}

All simulation and model fitting was performed in Python. Table~\ref{tab:ch4_stack}
lists the exact versions used; results were produced on a single machine
(Apple silicon, \texttt{arm64}, $12$ cores, $18$~GiB RAM) with no GPU and no
distributed compute.

\begin{table}[htbp]
\centering
\caption[Software versions]{Software versions used for all reported results.}
\label{tab:ch4_stack}
\begin{tabular}{ll}
\toprule
\textbf{Component} & \textbf{Version} \\
\midrule
Python            & 3.12.14 \\
PennyLane         & 0.45.1 \\
NumPy             & 2.5.3 \\
SciPy             & 1.18.1 \\
pandas            & 3.0.5 \\
scikit-learn      & 1.9.1 \\
Matplotlib        & 3.11.1 \\
Autograd          & bundled with PennyLane \\
\bottomrule
\end{tabular}
\end{table}

PennyLane provides both the noisy simulation and the quantum models, which removes
any question of consistency between the data-generating simulator and the simulator
used inside the hybrid arms. Two devices are used. The \texttt{default.mixed} device
performs density-matrix simulation and is used for all noisy circuit evaluation,
since the channels of Eqs.~\eqref{eq:ch3_depol} and~\eqref{eq:ch3_ad} are
non-unitary and cannot be represented on a state-vector simulator. The
\texttt{default.qubit} state-vector device is used inside the two hybrid arms, whose
embedding and ansatz circuits are unitary.

\section{Simulation procedure and cost}

For each circuit, the density matrix is evaluated once per noise scale
$\lambda \in \{0, 1, 2, 3\}$ and all five observables are then obtained from that
single matrix as $\mathrm{Tr}[O \rho]$, rather than by four separate circuit
executions per observable. This reduces the simulation cost of the primary dataset
by a factor of five and is possible only because the study is simulated; on hardware
each $(\text{circuit}, \text{observable}, \lambda)$ triple would be a separate
measurement job.

Generating the primary $900$-circuit dataset takes approximately three minutes at
$n = 5$ qubits. The five secondary datasets of $500$ circuits each take between
about one minute ($n = 4$) and about eight minutes ($n = 6$), the growth reflecting
the $4^n$ scaling of density-matrix simulation.

\section{Hyperparameter selection}

Hyperparameters are selected on the training fold only. For the gradient-boosted
trees and the quantum kernel the selection is four-fold grouped cross-validation,
with the circuit identifier as the grouping variable, so that no circuit appears in
both the inner-training and inner-validation parts of any fold. Two arms are the
exception. The ridge penalty is chosen by scikit-learn's efficient leave-one-out
cross-validation (\texttt{RidgeCV}), and the neural network stops early on a random
$10\%$ validation split of the training rows; neither of these inner splits is
grouped by circuit, so the selected values may be slightly optimistic. Every
\emph{test} evaluation remains grouped. Table~\ref{tab:ch4_grids}
lists the search spaces. No hyperparameter was selected using test-fold performance,
and the analytic arms have no hyperparameters to select.

\begin{table}[htbp]
\centering
\caption[Hyperparameter search spaces]{Hyperparameter search spaces. All searches use four-fold grouped
cross-validation on the training fold. The reduced grids in the right-hand column
were used for the sweep experiments of Sections~\ref{sec:ch5_noise},
\ref{sec:ch5_shots}, \ref{sec:ch5_qubits}, \ref{sec:ch5_learning},
\ref{sec:ch5_ablation} and~\ref{sec:ch5_transfer}, and for every realistic-regime
experiment of Sections~\ref{sec:ch5_scaling}--\ref{sec:ch5_drift}, to keep the total
runtime tractable. The main benchmark of Section~\ref{sec:ch5_main}, the
depth-extrapolation experiment, and the base-configuration points of the
noise-strength and system-size sweeps ($1.0\times$ and $n = 5$, which reuse the
main-benchmark fits) use the full grids and 40 epochs. The shot-limited kernel of
Section~\ref{sec:ch5_shotkernel} fixes $\beta = 0.15$ and selects $\alpha$ from 12
log-spaced points in $[10^{-6}, 10^{1}]$.}
\label{tab:ch4_grids}
\small
\begin{tabular}{p{3.4cm} p{5.3cm} p{3.6cm}}
\toprule
\textbf{Arm} & \textbf{Full grid} & \textbf{Reduced grid} \\
\midrule
Ridge regression &
$\alpha \in$ 30 points log-spaced in $[10^{-4}, 10^{3}]$ &
same \\
\addlinespace
Gradient-boosted trees &
leaves $\in \{15, 31\}$, learning rate $\in \{0.05, 0.1\}$, 300 iterations &
leaves $= 31$, rate $= 0.1$ \\
\addlinespace
Neural network &
fixed $(64, 64)$, early stopping, max 3000 iterations &
same \\
\addlinespace
Quantum kernel ridge &
bandwidth $\beta \in \{0.05, 0.15, 0.4, 1.0\}$, $\alpha \in$ 15 points
log-spaced in $[10^{-6}, 10^{1}]$ &
$\beta \in \{0.05, 0.15, 0.4\}$, $\alpha \in$ 8 points \\
\addlinespace
Data-reuploading VQC &
4 qubits, 4 layers, Adam, step $0.05$, batch $96$, 40 epochs &
30 epochs \\
\addlinespace
Bounded-encoding VQC &
identical to the above &
30 epochs \\
\bottomrule
\end{tabular}
\end{table}

The two variational arms are trained by minimizing mean squared error with the Adam
optimizer at a step size of $0.05$ and a batch size of $96$. Gradients are obtained
by automatic differentiation through the state-vector simulator rather than by the
parameter-shift rule; the two are mathematically equivalent for the expectation
values used here, and backpropagation is far cheaper in simulation. On hardware the
parameter-shift rule would be required, and the training cost would rise by a factor
of twice the parameter count per gradient evaluation -- a point returned to in
Section~\ref{sec:ch5_cost}.

\section{Random seeds and determinism}

Two seeds control all reported results. The dataset seed $20260912$ fixes the
circuit parameters and the binomial shot sampling; the model seed $7$ fixes the
train/test split, the cross-validation folds, the neural-network initialization,
the variational-circuit initialization and shuffling, and the bootstrap resampling.
Two families of repetition were added to put variance on the headline results
(Section~\ref{sec:ch5_robustness}): ten circuit-grouped splits with seeds
$100$--$109$, in which the scikit-learn model seeds follow the split; and ten
variational-circuit initialisation seeds $0$--$9$ on the reference split. Paired
bootstrap comparisons use $1000$ resamples of test circuits.

One implementation detail is worth recording because getting it wrong silently
invalidates the paired comparisons of Sections~\ref{sec:ch5_markovianity}
and~\ref{sec:ch5_drift}. The circuit parameters and the shot noise are drawn from
\emph{separate} generators, seeded at $20260912$ and $20260913$. A single shared
generator does not work: NumPy's binomial sampler uses rejection sampling, so the
number of underlying random words it consumes depends on its parameters, which in
turn depend on the noise configuration. With one stream, the circuit sequence
therefore becomes a function of the noise model, and the three realistic datasets --
which must differ only in their noise -- silently end up containing different
circuits. This was in fact observed during development: an initial generation run
produced datasets agreeing on only $185$ of $900$ circuit depths. Each dataset now
records a hash of its circuit sequence and of its ideal expectation values, and the
experiment driver refuses to run unless the hashes of all four realistic datasets
(production, drifted, defect-removed and reset control) agree.
Given these seeds, the entire pipeline is deterministic, with the single exception
noted below.

One implementation detail is worth recording because it affects reproducibility on
other machines. The scikit-learn grid searches are run single-threaded
(\texttt{n\_jobs=1}); process-based parallelism was unavailable in the execution
environment used here. Single-threaded execution removes any dependence of the
result on thread scheduling, so this constraint improves determinism at the cost of
wall-clock time.

\section{Source tree and file map}

Table~\ref{tab:ch4_files} maps each script to the outputs it produces. The pipeline
is designed so that the model code is independent of the data source: every script
downstream of data generation consumes a flat table with the columns
\texttt{circuit\_id}, \texttt{family}, \texttt{depth}, \texttt{n\_2q},
\texttt{observable}, \texttt{obs\_weight}, \texttt{noisy\_l1}, \texttt{noisy\_l2},
\texttt{noisy\_l3} and \texttt{ideal}. A table of hardware measurements with those
columns can therefore be substituted for the simulated one without modifying any
model code, which is the intended path for extending this work to a real device.

\begin{table}[htbp]
\centering
\caption[Source files and their outputs]{Source files and their outputs. Scripts are in \texttt{code/} and
outputs in \texttt{results/} of the accompanying project tree.}
\label{tab:ch4_files}
\small
\begin{tabular}{p{4.4cm} >{\raggedright\arraybackslash}p{9.0cm}}
\toprule
\textbf{Script} & \textbf{Produces} \\
\midrule
\texttt{gen\_data.py} &
Primary dataset (\texttt{benchmark\_data.csv}) and its metadata record. \\
\texttt{simulate.py} &
Secondary datasets for the noise-strength and qubit-count sweeps, parametrized by
qubit count, noise multiplier and shot budget. \\
\texttt{bench.py} &
Main benchmark: all nine arms on Splits A and B
(\texttt{benchmark\_results.csv}), per-observable and per-depth breakdown, and
the fitted predictions. \\
\texttt{experiments.py} &
The seven sweep experiments (\texttt{sweep\_results.csv}), including the
bounded-encoding arm. \\
\texttt{make\_figures.py} &
Every figure in Chapter~\ref{chap5:results}. \\
\texttt{simulate\_realistic.py} &
Device-realistic regime-II generator: inhomogeneous rates, coherent over-rotation,
$ZZ$ crosstalk, the non-Markovian defect, and both noise-scaling methods. Writes
\texttt{ds\_real\_A/B/M/R.csv} (production, drift, defect removed, defect reset)
with circuit and ideal-value signatures. \\
\addlinespace
\texttt{test\_folding.py} &
Noise-free unit test that unitary folding leaves every expectation value
unchanged, for both circuit families. \\
\addlinespace
\texttt{nonmarkov\_check.py} &
Breuer--Laine--Piilo measurement on the idling register
(\texttt{nonmarkovianity\_check.json}) and the coupling sweep. \\
\addlinespace
\texttt{experiments\_realistic.py} &
Regime-II experiments: scaling-method comparison, both memory controls, calibration
drift, shot-limited quantum kernel with five draws per budget
(\texttt{sweep\_realistic.csv}) and all paired comparisons
(\texttt{paired\_realistic.csv}). \\
\addlinespace
\texttt{stats\_util.py} &
Paired, circuit-grouped bootstrap on RMSE differences. \\
\texttt{robustness.py} &
Ten circuit-grouped splits per configuration with paired comparisons
(\texttt{robust\_splits\_*.csv}, \texttt{robust\_pairs\_*.csv}). \\
\texttt{followups.py} &
Depth extrapolation by input set and learning-curve variance over $20$ draws
(\texttt{followup\_results.csv}, \texttt{followup\_paired.csv}). \\
\texttt{vqc\_seeds.py} &
Variational arms over ten initialisation seeds (\texttt{vqc\_seed\_sweep.csv}). \\
\texttt{make\_appendixB.py} &
Regenerates the regime-II, robustness and non-Markovianity tables of
Appendix~\ref{AppendixB} from the result files. \\
\texttt{make\_figures\_realistic.py} &
The regime-II and robustness figures of Chapter~\ref{chap5:results}. \\
\addlinespace
\texttt{qec\_gen\_data.py} &
Repetition-code syndrome datasets under phenomenological noise
(Section~\ref{chap6:decoding} states why this substitutes for the surface code). \\
\texttt{qec\_bench.py} &
Decoder benchmark (\texttt{qec\_decoder\_results.csv}) and the figures of
Chapter~\ref{chap6:decoding}. \\
\bottomrule
\end{tabular}
\end{table}

\section{Verification practices}\label{sec:ch4_verification}

Five practices are worth stating because they affected the reported numbers. First,
every figure caption and every claim about a trend was checked against the saved
result table before being written, and one panel title was corrected during
preparation when the table showed the trend it asserted did not hold for all
plotted series. Second, the quantum-kernel arm was initially run at unit bandwidth
and produced an error worse than the unmitigated baseline; rather than report that
as the performance of quantum kernel methods, bandwidth was promoted to a tuned
hyperparameter, and the untuned result is reported in
Section~\ref{sec:ch5_bandwidth} as a sensitivity finding in its own right. Both
episodes are recorded here because they illustrate a general hazard in this
literature: a hybrid model reported without its critical hyperparameter tuned is
evidence about the implementation, not about the method.

Third, every circuit transformation used for noise scaling is unit-tested in the
absence of noise. Unitary folding must leave every expectation value unchanged
($L L^{\dagger} L = L$), and \texttt{test\_folding.py} checks this for both circuit
families. The test was written after a pre-submission review found that the first
version of the realistic simulator built the inverse of a hardware-efficient layer
with its two rotations in forward order ($R_Y^{\dagger}$ before $R_Z^{\dagger}$
instead of after), so the folded circuits computed a different unitary and every
folded feature of that family was invalid. After the fix the folded and unfolded
values agree to $3.9 \times 10^{-16}$, the four realistic datasets were regenerated,
and every regime-II result in Chapter~\ref{chap5:results} comes from the corrected
data. The parameter-scaled columns were unaffected.

Fourth, comparisons between two methods are made with a paired, circuit-grouped
bootstrap on the difference of their errors rather than by inspecting whether two
marginal intervals overlap, and every headline comparison is repeated over ten
splits; Section~\ref{sec:ch5_robustness} reports what changed as a result.

Fifth, every quantitative statement in the abstract, the contribution list, the
chapter summaries and the conclusions was audited against the saved result tables
before submission; statements that the tables did not support were corrected, and
the audit log is part of the accompanying project tree
(\texttt{results/claims\_audit.csv}).

\end{document}
